\documentclass[10pt,a4paper]{amsart}
\usepackage[margin=19mm]{geometry}
\usepackage{amsmath,amssymb,mathtools}
\usepackage[scr=rsfs,cal=euler]{mathalfa}
\usepackage{xfrac}
\usepackage[unicode,hidelinks,bookmarks=false]{hyperref}
\newcommand{\ie}{i.e.\ }
\newcommand{\cf}{cf.\ }
\newcommand{\resp}{respectively\ }
\newcommand{\Subsection}{\subsection}
\providecommand{\nobreakdash}{\nobreak\hbox{-}}
\setlength{\emergencystretch}{2em}
\multlinegap0pt
\usepackage{amssymb}
\usepackage{enumerate}
\usepackage{dsfont}
\usepackage{tikz}
\newtheorem{theorem}{Theorem}
\newtheorem{lemma}{Lemma}
\usepackage{cleveref}
\crefname{theorem}{Theorem}{Theorems}
\crefname{lemma}{Lemma}{Lemmas}
\newcommand{\DAS}[1]{{\color{red} #1}}
\newcommand{\E}{\mathbb{E}}
\def\MRV{\mathcal{MRV}}
\renewcommand\P{\mathbb{P}}
\newcommand{\R}{\mathbb{R}}
\newcommand{\VF}[1]{{\color{blue} #1}}
\def\bZ{\boldsymbol Z}
\newcommand{\beam}{\begin{eqnarray}}
\newcommand{\beao}{\begin{eqnarray*}}
\def\bz{\boldsymbol z}
\def\cR{\mathcal{R}}
\newcommand{\eeam}{\end{eqnarray}\noindent}
\newcommand{\eeao}{\end{eqnarray*}\noindent}
\title{Aggregating heavy-tailed random vectors:\\ from finite sums to L\'evy processes}
\author{Bikramjit Das, Vicky Fasen-Hartmann}
\date{Source: arXiv:2301.10423v1 (CC BY 4.0)}
\begin{document}
\maketitle

\noindent\emph{Source and licence.} Aggregating heavy-tailed random vectors:\\ from finite sums to L\'evy processes. Version arXiv:2301.10423v1 (CC BY 4.0), distributed by arXiv under the Creative Commons Attribution 4.0 International licence, http://creativecommons.org/licenses/by/4.0/, as recorded in the arXiv licence metadata for this exact version and shown on the abstract page https://arxiv.org/abs/2301.10423v1. Full licence notice: \texttt{licenses/arxiv-2301.10423-CC-BY-4.0.txt}.\par\smallskip

\noindent\textit{Editorial scope.} Counted unit: the article's lemma 3.4 (source0.tex lines 1667--1688). The article's complete original proof of it is delivered with it (source0.tex lines 1691--1782, one \texttt{proof} environment of 92 lines). No mathematics is written, repaired or re-derived: the statement and the proof are the article's own lines, in the article's own notation, and no retained line is substituted or re-flowed.  Retained uncounted as the source-local context the proof needs: lines 808--838; lines 1633--1640.  Nothing was omitted from any retained span, so the package carries no omission record.  No retained line contains a citation, so the wrapper carries no bibliography and no bibliography entry.  No retained line contains a figure environment, an image-inclusion command or drawing code, so no image is delivered and the package declares no asset dependency.  Editorial formatting only, in addition: an AMS \texttt{amsart} preamble that restates the article's own declarations and macros used by the retained lines, namely the environments \texttt{theorem}, \texttt{lemma} on separate counters, together with the standard packages those lines need.  The retained environments print in this document as Theorem 1, Lemma 1 in order of appearance, whereas the article prints the same items with the numbering of its own full document.  The complete native source of the article as distributed in the arXiv e-print for this exact version is bundled unchanged as \texttt{sources/136497/source0.tex}.
\begin{theorem} \label{mainTheorem}
Let $\bZ^{(1)}, \bZ^{(2)} \in \R_+^d$ be independent random vectors, each with tail equivalent marginal distributions and $\{\bZ^{(k)}\in \MRV^*(\alpha_i^{(k)},b_i^{(k)},\mu_i^{(k)},\E_d^{(i)}), i=1,\ldots,d; \Delta_k\}$ for $k=1,2$, i.e., they are adapted-MRV on $\R^d_{+}$. Define $\alpha_0^{(k)}=0, \, b_0^{(k)\leftarrow}(t)\equiv 1,\, \mu_0^{(k)}\equiv 1$ and
\[ I(i): =   \,\mathrm{argmax}_{j \in \{0,\ldots,i\}}  \{\bar{c}_j^{(i)}: \bar{c}_j^{(i)}<\infty\}\]
where
\[ \bar{c}_j^{(i)}:= \max_{0\le m\le i}\left\{\limsup_{t\to\infty} \frac{b_{j}^{(1)\leftarrow}(t)b_{i-j}^{(2)\leftarrow}(t)}{b_{m}^{(1)\leftarrow}(t)b_{i-m}^{(2)\leftarrow}(t)} \right\}.\]
Define for $i=1,\ldots, d$, and $m=0,\ldots, i$:
$$c_m^{I(i)}:=\lim_{t\to\infty}\frac{b_{I(i)}^{(1)\leftarrow}(t)b_{i-I(i)}^{(2)\leftarrow}(t)}{b_{m}^{(1)\leftarrow}(t)b_{i-m}^{(2)\leftarrow}(t)}.$$
%\frac{b_{i^*}^{(1)\leftarrow}(t)b_{i-i^*}^{(2)\leftarrow}(t)}{b_{m}^{(1)\leftarrow}(t)b_{i-m}^{(2)\leftarrow}(t)} \right\} \]
%
% There exists an $i^*\in\{0,\ldots,i\}$ with $$\limsup_{t\to\infty}\frac{b_{i^*}^{(1)\leftarrow}(t)b_{i-i^*}^{(2)\leftarrow}(t)}{b_{m}^{(1)\leftarrow}(t)b_{i-m}^{(2)\leftarrow}(t)}<\infty \quad \text{ for any }  \quad m\in\{0,\ldots,i\}$$
% where $b_0^{(1)\leftarrow}(t):=b_0^{(2)\leftarrow}(t):=1$ and set
%$$c_m^{i^*}=\lim_{t\to\infty}\frac{b_{i^*}^{(1)\leftarrow}(t)b_{i-i^*}^{(2)\leftarrow}(t)}{b_{m}^{(1)\leftarrow}(t)b_{m-i^*}^{(2)\leftarrow}(t)} \quad \text{ for } \quad m=0,\ldots,i.$$  Define $$M(i^*):=\left\{m\in\{0,\ldots,i\}:\limsup_{t\to\infty}\frac{b_{i^*}^{(1)\leftarrow}(t)b_{i-i^*}^{(2)\leftarrow}(t)}{b_{m}^{(1)\leftarrow}(t)b_{m-i^*}^{(2)\leftarrow}(t)}>0\right\}.$$
Suppose that for  $k=1,2, $ and each $m=1,\ldots,i-1$, either
$c_m^{I(i)}=0$ or  $\lim_{t\to\infty}b_m^{(k)\leftarrow}(t)/b_{m+1}^{(k)\leftarrow}(t)=0$.
 Then
 \beao
    \{\bZ^{(1)}+\bZ^{(2)}\in\MRV^*(\alpha_{i},b_{i},\mu_{i}^{\oplus},\E_d^{(i)}), i=1,\ldots,d;\Delta^{\oplus}\}
 \eeao
% \marginpar{\VF{What is $\Delta^{\oplus}$?}\DAS{good question :))}}
 with $\alpha_i=  \alpha_{I(i)}^{(1)} +\alpha_{i-I(i)}^{(2)}, b_{i}^{\leftarrow}(t)=b_{I(i)}^{(1)\leftarrow}(t)b_{i-{I(i)}}^{(2)\leftarrow}(t), $ and
 \beao
    \mu_i^{\oplus}(A)=\sum_{m=0}^i c_m^{I(i)}\mu_{m,i}^*(A) \quad \text{ for }A\in\mathcal{B}(\E_d^{(i)}),
 \eeao
 where $\mu_{m,i}^*$ is the measure which is uniquely defined on $\cR^{(i)}$ as follows: for \linebreak $A=\{\bz \in \R_+^d: z_j>x_j \;\forall\; j\in S\}\in \cR^{(i)}$ with $S\subseteq \mathbb{I}$, $|S|\ge i$ and $x_j>0$   $\forall j\in S$ we have % $\mu_{m,i}^*(A):=0$ if $\lim_{t\to\infty}b_i^{\leftarrow}(t)/b_{|S|}^{\leftarrow}(t)=0$ and else
 \begin{align*}
        &\mu_{m,i}^*\left(A\right)\\
         &=\sum_{\genfrac{}{}{0pt}{}{J\subseteq S\cup \{\emptyset\}}{|J|=m}} \mu_{m}^{(1)}\Bigg(\{\bz\in\E_d^{(m)}: z_{j}>x_{j}\;\forall\;j\in J\}\Bigg)
        \mu_{i-m}^{(2)}\,\Bigg(\{\bz\in\E_d^{(i-m)}:  z_{j}>x_{j}\;\forall\; j\in S\backslash J\}\Bigg).
 \end{align*}
 Moreover, $\Delta^{\oplus} \in \{\max(\Delta_1+1,\Delta_2+1), \ldots, \min(\Delta_1+\Delta_2,d)\}$.
 \end{theorem}

\begin{lemma} \label{Lemma 4.2}
Let the assumptions of \Cref{mainTheorem} hold. Then for any
$m=0,\ldots,i-1$:
\begin{align}
    \lim_{t\to\infty}b_{I(i)}^{(1)\leftarrow}(t)b_{i-I(i)}^{(2)\leftarrow}(t)\,\P(Z_{(m+1)}^{(1)}>t)\,\P(Z_{(i-m)}^{(2)}>t)=& 0 \label{appeq:B1}\\ \intertext{ and }
   { \lim_{t\to\infty}b_{I(i)}^{(1)\leftarrow}(t)b_{i-I(i)}^{(2)\leftarrow}(t)\,\P(Z_{(m+1)}^{(2)}>t)\,\P(Z_{(i-m)}^{(1)}>t)}=& 0. \label{appeq:B2}
\end{align}
\end{lemma}

\begin{lemma} \label{Lemma 3.4}
Let the assumptions of \Cref{mainTheorem} hold and
%Let $\bZ^{(1)}$ and $\bZ^{(2)}$ be independent random vectors in $\R_+^d$ with identically distributed marginal distributions and let %$\bZ^{(k)}\in\MRV(\alpha_i^{(k)},b_i^{(k)},\mu_i^{(k)},\E_d^{(i)})$  for $i=1,\ldots,d$ and $k=1,2$. Define $\alpha_0^{(k)}=0, %b_0^{(k)\leftarrow}(t)\equiv 1, k=1,2$, and
%\[ I(i): =  \arg \max_{j \in \{0,\ldots,i\}}  \{\bar{c}_j: \bar{c}_j<\infty\}\]
%where
%\[ \bar{c}_j:= \max_{1\le m\le i}\left\{\limsup_{t\to\infty} %\frac{b_{j}^{(1)\leftarrow}(t)b_{i-j}^{(2)\leftarrow}(t)}{b_{m}^{(1)\leftarrow}(t)b_{i-m}^{(2)\leftarrow}(t)} \right\}.\]
%Also define for $i=1,\ldots, d$, and $m=0,\ldots, i,$:
%$$c_m^{I(i)}=\lim_{t\to\infty}\frac{b_{I(i)}^{(1)\leftarrow}(t)b_{i-I(i)}^{(2)\leftarrow}(t)}{b_{m}^{(1)\leftarrow}(t)b_{i-m}^{(2)\leftarrow}(t)}.$$
%\frac{b_{i^*}^{(1)\leftarrow}(t)b_{i-i^*}^{(2)\leftarrow}(t)}{b_{m}^{(1)\leftarrow}(t)b_{i-m}^{(2)\leftarrow}(t)} \right\} \]
%
% There exists an $i^*\in\{0,\ldots,i\}$ with $$\limsup_{t\to\infty}\frac{b_{i^*}^{(1)\leftarrow}(t)b_{i-i^*}^{(2)\leftarrow}(t)}{b_{m}^{(1)\leftarrow}(t)b_{i-m}^{(2)\leftarrow}(t)}<\infty \quad \text{ for any }  \quad m\in\{0,\ldots,i\}$$
% where $b_0^{(1)\leftarrow}(t):=b_0^{(2)\leftarrow}(t):=1$ and set
%$$c_m^{i^*}=\lim_{t\to\infty}\frac{b_{i^*}^{(1)\leftarrow}(t)b_{i-i^*}^{(2)\leftarrow}(t)}{b_{m}^{(1)\leftarrow}(t)b_{m-i^*}^{(2)\leftarrow}(t)} \quad \text{ for } \quad m=0,\ldots,i.$$  Define $$M(i^*):=\left\{m\in\{0,\ldots,i\}:\limsup_{t\to\infty}\frac{b_{i^*}^{(1)\leftarrow}(t)b_{i-i^*}^{(2)\leftarrow}(t)}{b_{m}^{(1)\leftarrow}(t)b_{m-i^*}^{(2)\leftarrow}(t)}>0\right\}.$$
 $A= \{\bz \in \R_+^d: z_j>x_j \;\forall\; j\in S\}$ be a rectangular set with
 $S\subseteq \mathbb{I}$, $|S|= i$, $x_{j}>0$ for $j\in S$ and $\mu_i^{\oplus}(\partial A)=0$. Then
\beam\label{eq:lemma2}
   \lim_{t\to\infty}b_{I(i)}^{(1)\leftarrow}(t)b_{i-I(i)}^{(2)\leftarrow}(t) \P\left(\bZ^{(1)}+\bZ^{(2)}\in tA\right)
  % &&=\sum_{m=0}^ic_{m}^{I(i)}\sum_{J\subseteq S\cup \{\emptyset\} \atop |J|=m} \mu_{m}^{(1)}\left(\left\{\bz\in\E_d^{(m)}: \bigcap_{j\in J} \{z_{j}>x_{j}\}\right\}\right)
  %      \mu_{i-m}^{(2)}\left(\left\{\bz\in\E_d^{(i-m)}: \bigcap_{j\in S\backslash J}\{ z_{j}>x_{j}\}\right\}\right)\\
        =\mu_i^{\oplus}(A).
\eeam
\end{lemma}

\begin{proof} $\mbox{}$\\
 \textbf{Step 1.} First, we derive an upper bound for the left hand side of \eqref{eq:lemma2}.  Let $0<\epsilon<1$
 such that $\epsilon x_j<(1-\epsilon)x_l$ for all $j,l\in S$.
 Then
\beam \label{A6}
    \P\left(\bZ^{(1)}+\bZ^{(2)}\in tA\right)&=&\Bigg[\sum_{\genfrac{}{}{0pt}{}{J_2\subseteq S\cup \{\emptyset\}}{J_1=S\backslash J_2}}+\sum_{\genfrac{}{}{0pt}{}{J_2\subseteq  S\cup\{\emptyset\}}{J_1 \subsetneq S\backslash J_2\cup \{\emptyset\}}}\Bigg]\P\left(\bigcap_{j\in S}\left\{Z_j^{(1)}+Z_j^{(2)}>tx_j\right\}
                 \right. \nonumber\\
        &&\quad\quad \cap\bigcap_{j\in J_1}\left\{Z_j^{(1)}\leq t\epsilon x_j\right\}\cap\bigcap_{j\in S\backslash J_1}\left\{Z_j^{(1)}> {t (1-\epsilon)} x_j\right\} \nonumber\\
        &&\quad\quad\left. \cap\bigcap_{j\in J_2}\left\{Z_j^{(2)}\leq t\epsilon x_j\right\}\cap\bigcap_{j\in S\backslash J_2}\left\{Z_j^{(2)}> {t(1-\epsilon)} x_j\right\} \right) \nonumber\\
        &&=:M_{1}(t,\epsilon)+M_{2}(t,\epsilon).
\eeam
Note that in case $J_1\cap J_2 \not=\emptyset$ these probabilities are zero since it results in computing probabilities of empty sets.
Next, we  find an upper bound for $M_{1}(t,\epsilon)$. Note that
\beao
    M_{1}(t,\epsilon)&\leq&\sum_{\genfrac{}{}{0pt}{}{J_2\subseteq S \cup\{\emptyset\}}{J_1=S\backslash J_2}}
        \P\left(\bigcap_{j\in S\setminus J_1}\left\{Z_j^{(1)}> t(1-\epsilon) x_j\right\}\cap\bigcap_{j\in S \setminus J_2}\left\{Z_j^{(2)}> t(1-\epsilon) x_j\right\}\right)\\
        &=&\sum_{J_2\subseteq S\cup\{\emptyset\} }
       \P\left(\bigcap_{j\in J_2}\left\{Z_j^{(1)}> t(1-\epsilon) x_j\right\}\right)
        \P\left(\bigcap_{j\in S\backslash J_2}\left\{Z_j^{(2)}> t(1-\epsilon) x_j\right\}\right).
\eeao
Let $A_{\epsilon}:= \{\bz \in \R_+^d: z_j>(1-\epsilon)x_j \,\forall\, j\in S\}$ and choose $\epsilon>0$ such that $\mu_i^{\oplus}(\partial A_{\epsilon})=0$.
Then
\beam \label{A5}
    %\lefteqn{
    \limsup_{t\to\infty}b_{I(i)}^{(1)\leftarrow}(t)b_{i-I(i)}^{(2)\leftarrow}(t)M_{1}(t,\epsilon)
    \leq \mu_i^{\oplus}(A_{\epsilon}). %\sum_{m=0}^i c_{m}^{I(i)}\sum_{J\subseteq S\cup \{\emptyset\} \atop |J|=m} \mu_{m}^{(1)}\left(\left\{\bz\in\E_d^{(m)}: \bigcap_{j\in J} \{z_{j}>x_{j}(1-\epsilon)\}\right\}\right)
        %\mu_{i-m}^{(2)}\left(\left\{\bz\in\E_d^{(i-m)}: \bigcap_{j\in S\backslash J}\{ z_{j}>x_{j}(1-\epsilon)\}\right\}\right). \nonumber
\eeam
Define $x^*:=\min_{j\in S} x_j$. Following a similar argument for $M_{2}(t,\epsilon)$ we get the upper bound
\beao
    M_{2}(t,\epsilon)&\leq& \sum_{\genfrac{}{}{0pt}{}{J_2\subseteq S\cup\{\emptyset\}}{ J_1\subsetneq S\backslash J_2\cup\{\emptyset\}}}
        \P\Big(\bigcap_{j\in J_2\cup S\backslash J_1}\left\{Z_j^{(1)}> t(1-\epsilon) x_j\right\}\Big)\P\Big(\bigcap_{j\in  J_1\cup S\backslash J_2}\left\{Z_j^{(2)}> t(1-\epsilon) x_j\right\}\Big)\\
    &\leq& \sum_{\genfrac{}{}{0pt}{}{J_2\subseteq S\cup\{\emptyset\}}{ J_1\subsetneq S\backslash J_2\cup\{\emptyset\}}}
        \P\left(Z_{(|J_2\cup S\backslash (J_1\cup J_2)|)}^{(1)}> t(1-\epsilon) x^*\right)\,\P\left(Z_{(|S\backslash J_2|)}^{(2)}> t(1-\epsilon) x^*\right)\\
    &\leq &\sum_{m=0}^{i-1}\sum_{l=1}^{i-m}\sum_{\genfrac{}{}{0pt}{}{J_2\subseteq S\cup\{\emptyset\} }{ |J_2|=m }}\sum_{\genfrac{}{}{0pt}{}{J_1\subsetneq S\backslash J_2\cup\{\emptyset\}}{|S\backslash ( J_1 \cup J_2)|=l}}
        \P\left(Z_{(m+l)}^{(1)}> t(1-\epsilon) x^*\right)\,\P\left(Z_{(i-m)}^{(2)}> t(1-\epsilon) x^*\right).
\eeao
Finally, an application of \Cref{Lemma 4.2} yields
\beam
    \lefteqn{\limsup_{t\to\infty}b_{I(i)}^{(1)\leftarrow}(t)b_{i-I(i)}^{(2)\leftarrow}(t) M_{2}(t,\epsilon)}\nonumber\\
    && \quad\quad\leq \sum_{m=0}^{i-1}\sum_{l=1}^{i-m}\sum_{\genfrac{}{}{0pt}{}{J_2\subseteq S\cup\{\emptyset\}}{ |J_2|=m }}\sum_{\genfrac{}{}{0pt}{}{J_1\subsetneq S\backslash J_2 \cup\{\emptyset\}}{|S\backslash ( J_1 \cup J_2)|=l}}{\limsup_{t\to\infty} }\, b_{I(i)}^{(1) \leftarrow}(t)b_{i-I(i)}^{(2)\leftarrow}(t) \nonumber\\
    && \quad\quad\quad\quad\quad \P\left(Z_{(m+1)}^{(1)}> t(1-\epsilon) x^*\right)\,\P\left(Z_{(i-m)}^{(2)}> t(1-\epsilon) x^*\right)  =0. \label{A4}
\eeam
Now from \eqref{A6}, \eqref{A5} and \eqref{A4} we have
\beao
     \limsup_{t\to\infty}b_{I(i)}^{(1)\leftarrow}(t)b_{i-I(i)}^{(2)\leftarrow}(t)\,\P\left(\bZ^{(1)}+\bZ^{(2)}\in tA\right)
        \leq \mu_i^{\oplus}(A_{\epsilon})\downarrow \mu_i^{\oplus}(A) \quad \text{ as }\epsilon\downarrow 0, %\sum_{m=0}^i c_{m}^{I(i)}\sum_{J\subseteq S\cup \{\emptyset\} \atop |J|=m} \mu_{m}^{(1)}\left(\left\{\bz\in\E_d^{(m)}: \bigcap_{j\in J} \{z_{j}>x_{j}\}\right\}\right)
        %\mu_{i-m}^{(2)}\left(\left\{\bz\in\E_d^{(i-m)}: \bigcap_{j\in S\backslash J}\{ z_{j}>x_{j}\}\right\}\right).
\eeao
where in the last step we use the fact that $\mu_i^{\oplus}(\partial A)=0$.

 \textbf{Step 2.} Next, we derive a lower bound for the asymptotic limit.
There are a total of $2^{|S|}$ subsets of $S$ which we order as
$J(1),\ldots, J(2^{|S|})$ (in any way). Now, define the sets
\beao
    C_{l}:=\bigcap_{j\in J(l)}\left\{Z_j^{(1)}> t x_j\right\}\cap \bigcap_{j\in S\backslash J(l)}\left\{Z_j^{(2)}> t x_j\right\}, \quad
        l=1,\ldots,2^{|S|}.
\eeao
Then,
\beao
    \left\{\bZ^{(1)}+\bZ^{(2)}\in tA\right\}\supseteq\bigcup_{l=1}^{2^{|S|}}C_l,
\eeao
and using the inclusion-exclusion principle we have
\beam \label{A1}
    \P\left(\bZ^{(1)}+\bZ^{(2)}\in tA\right)\geq \P\left(\bigcup_{l=1}^{2^{|S|}}C_l\right)% \nonumber\\
        \geq \sum_{l=1}^{2^{|S|}}\P(C_l)-\sum_{1\leq l_1<l_2\leq 2^{|S|}}\P(C_{l_1}\cap C_{l_2}).
\eeam
Now on one hand,
\beam \label{A2}
    \lim_{t\to\infty}b_{I(i)}^{(1)\leftarrow}(t)b_{i-I(i)}^{(2)\leftarrow}(t)\sum_{l=1}^{2^{|S|}}\P(C_l)=\mu_i^{\oplus}(A),
    %    &&=\sum_{m=0}^i c_{m}^{I(i)}\sum_{J\subseteq S\cup \{\emptyset\} \atop |J|=m} \mu_{m}^{(1)}\left(\left\{\bz\in\E_d^{(m)}: \bigcap_{j\in J} \{z_{j}>x_{j}\}\right\}\right)
     %   \mu_{i-m}^{(2)}\left(\left\{\bz\in\E_d^{(i-m)}: \bigcap_{j\in S\backslash J}\{ z_{j}>x_{j}\}\right\}\right).  \nonumber
\eeam
and on the other hand, for any $1\leq l_1<l_2\leq 2^{|S|}$ the inequality
\beao
    0\leq \P(C_{l_1}\cap C_{l_2})\leq \P(Z_{(|J(l_1)\cup J(l_2)|)}^{(1)}>t x^*)\P(Z_{(|S\backslash (J(l_1)\cap J(l_2))|) }^{(2)}>tx^*)
\eeao
holds.  Define $m:=|J(l_1)\cap J(l_2)|$. Since $J(l_1)\not=J(l_2)$ we have $|J(l_1)\cup J(l_2)|\geq |J(l_1)\cap J(l_2)|+1=m+1$.
Hence, a conclusion of \Cref{Lemma 4.2} is that
\beam \label{A3}
    \lefteqn{\limsup_{t\to\infty}b_{I(i)}^{(1)\leftarrow}(t)b_{i-I(i)}^{(2)\leftarrow}(t)\,\P(C_{l_1}\cap C_{l_2})}  \\
        &&\leq \limsup_{t\to\infty}b_{I(i)}^{(1)\leftarrow}(t)b_{i-I(i)}^{(2)\leftarrow}(t)
         \P(Z_{(m+1)}^{(1)}>t x^*)\P(Z_{(i-m) }^{(2)}>tx^*) =0 \nonumber
\eeam
for any $1\leq l_1<l_2\leq 2^{|S|}$. Then \eqref{A1}, \eqref{A2} and \eqref{A3} result in the lower bound
\beao
    \liminf_{t\to\infty}b_{I(i)}^{(1)\leftarrow}(t)b_{i-I(i)}^{(2)\leftarrow}(t)\,\P\left(\bZ^{(1)}+\bZ^{(2)}\in tA\right)
    \geq \mu_i^{\oplus}(A).%\sum_{m=0}^i c_{m}^{I(i)}\sum_{J\subseteq S\cup \{\emptyset\} \atop |J|=m} \mu_{m}^{(1)}\left(\left\{\bz\in\E_d^{(m)}: \bigcap_{j\in J} \{z_{j}>x_{j}\}\right\}\right)
%        \mu_{i-m}^{(2)}\left(\left\{\bz\in\E_d^{(i-m)}: \bigcap_{j\in S\backslash J}\{ z_{j}>x_{j}\}\right\}\right)
\eeao
and together with the upper bound in Step~1 the lemma is proven.
\end{proof}

\end{document}
